\section{统一的动作级建模与调度}

本节先在 \S\ref{sec:design_modeling} 中给出动作级统一建模，覆盖异构外部资源与多样化工具类型；然后在 \S\ref{sec:design_scheduling} 中介绍我们的弹性资源调度算法。

\subsection{动作建模}
\label{sec:design_modeling}

\paraf{向量化资源成本建模。}
为了实现统一的动作级调度，首先需要对资源成本进行清晰而统一的表示。给定一个动作 $a_i$，它可能同时消耗多种资源类型。例如，DeepSearch 任务中的一个动作可能会并行或串行地访问多个网站，而每个网站都有独立的 API 配额，因此可以被视作不同的资源类型；AI 编程任务中的一个动作则既依赖 CPU 也依赖内存。因此，我们将动作成本表示为多维向量 $C_i = (c_{i,0}, c_{i,1},...,c_{i,k-1})$，其中 $k$ 表示由 \sysname 管理的资源类型数量。该向量使调度器可以联合考虑多种资源约束。向量的每个维度分别对应一种资源消耗，包括 CPU、GPU、内存、显存、网络带宽以及外部 API 配额等。

更重要的是，单个资源维度上的成本并不是静态值。对于对资源数量敏感的弹性动作，其成本通常表示为从最小可行值到最大可行值的范围。更复杂的情况是，一些多 GPU 动作只允许离散资源单位，例如 1、2、4、8。因此，$C_i$ 中的每个 $c_{i,j}$ 都附带自己的可行资源集合。这种建模方式使调度器能够在不破坏可行性的前提下灵活调整资源分配，并在调度时利用先验缩小搜索空间。

\parabf{弹性建模。}
对弹性动作，我们额外引入参数来描述其弹性行为。对于某一特定资源，其弹性被建模为：当分配 $m$ 个资源单位时，动作可获得一个介于 0 和 1 之间的弹性比例。文中采用如下形式：
\begin{equation}
\label{equ:scale}
\begin{aligned}
a.\text{getDur}(m) = \frac{T_{ori}}{E(m) * m}\quad (0<E(m)\le1),
\end{aligned}
\end{equation}
其中 $T_{ori}$ 表示动作在单个资源单位下的执行时长。为了简化建模与优化，我们假设一个动作的所有资源中，只有一种资源类型主导其弹性。例如，对 GPU 动作而言，显存是执行必要条件，但真正决定时延的通常是并行度。因此，动作的弹性可以视作关于关键资源类型使用量的映射。

\parabf{执行时长建模。}
为了服务于后续的弹性调度算法，动作建模中除了弹性以外，还需要显式纳入执行时长。因此，我们为可剖析的动作记录其原始执行时长，并按单个资源单位归一化。配合弹性定义，系统便能够估计不同资源分配下的时间代价，这对优化调度决策、在性能与资源效率之间做平衡是必要的。

\subsection{弹性资源调度}
\label{sec:design_scheduling}

在动作级资源分配与调度的设定下，我们需要一个面向有限资源条件的弹性调度算法来保证效率。我们把动作完成时间（ACT）作为优化目标。类似于作业完成时间（JCT），ACT 可分解为两部分：(1) 动作排队等待被调度的时间；(2) 动作实际执行时间。因此，调度目标可以写成最小化所有动作 ACT 总和：
\begin{equation}
\label{equ:obj}
\begin{aligned}
 \text{Objective:} \quad {\text{minimize}} ~ ACTs =\sum_{0\leq i < N} (T^q_{i} + T_i),
\end{aligned}
\end{equation}
其中 $N$ 表示调度时等待队列中的动作数，$T_i$ 为第 $i$ 个动作的实际执行时长，$T^q_i$ 为其排队时间。

具体到调度问题，它包含两个相互耦合的子问题：(1) 确定等待队列中动作的执行顺序；(2) 确定每个动作应分配多少资源单位。这两个问题都不容易。为了缩小搜索空间，并围绕本文关注的核心现象“资源过量预留”展开，我们主要优化第二个问题，同时对第一个问题采用简单但稳定的 First-Come First-Served（FCFS）策略。由于所有动作都位于各自轨迹的关键路径上，智能体 RL 对 starvation 十分敏感，因此 FCFS 是一个稳妥选择。当然，我们的调度框架也可以与其他排序策略结合。

第二个问题本质上是一个经典难题：在有限资源下为多个弹性作业分配资源，并最小化总 JCT，该问题已知是 NP-hard~\citep{bruno1974scheduling}。又由于调度窗口极短，同时资源类型和作业模式都高度异构，因此我们采用启发式方法。总体上，我们会先根据资源最小需求选出一组候选动作，然后按其关键弹性资源分组，再分别执行弹性分配。对于可缩放动作，我们使用带目标近似的贪心驱逐算法：先用最小资源计算初始目标值，然后逐步移除候选集合尾部动作，并把其资源回收给前面的动作；若目标值下降，则继续驱逐，否则停止。为估计目标，我们把候选集合的 ACT 精确计算，而把等待队列剩余动作的 ACT 通过近似方式估计，从而控制调度开销。

\begin{algorithm}[t]
\caption{弹性资源调度算法。}\label{alg:ersa}
\begin{footnotesize}
\begin{algorithmic}[1]
\State 候选动作先按照 FCFS 从等待队列中选出
\State 每个动作先分配其最小可行资源
\State 按关键弹性资源将动作分组
\State 对每组动作分别执行弹性分配
\State 通过近似目标函数比较不同驱逐与分配方案
\State 输出本轮真正进入执行的动作及其资源分配
\end{algorithmic}
\end{footnotesize}
\end{algorithm}

对于既可伸缩又已知弹性与时长信息的动作，我们提出了基于目标近似的贪心算法。它先假设所有候选动作均采用最小资源，再不断尝试驱逐队尾动作并把资源重新分给保留下来的动作；一旦继续驱逐不再降低目标值，算法便停止。候选集合内部的资源分配通过一个拓扑无关的动态规划（DP）算法求解，而等待队列剩余动作则通过顺序插入完成堆的方式进行近似。为了进一步提高近似精度，系统还引入了参数 $depth$，允许最靠前的剩余动作探索更多资源分配选项。实践中，$depth=2$ 或 $3$ 已经足够。

对不可伸缩资源或弹性未知的动作，算法则直接选择这些候选动作，并仅为其分配最小所需资源，同时保证所有向量化资源约束满足。

\section{异构资源管理器}
\label{sec:design_resource_managing}

虽然我们给出了统一抽象和调度算法，但要真正实现 \textbf{Breakdown} 与 \textbf{Pool}（\S\ref{sec:background_opp}），仍然需要资源特定的管理策略。具体而言，\textbf{Breakdown} 关注如何在释放资源的同时保留执行状态或支持快速恢复，也就是高效上下文切换；同时还要支持每个动作的弹性并行度。\textbf{Pool} 则关注如何设计资源分配方式以缓解碎片化并提高并行效率，从而为统一的弹性调度框架提供支撑。正如 \S\ref{sec:background_opp} 所述，不同资源在特征和拓扑结构上差异很大，因此我们为不同资源类型设计了不同管理器。

\subsection{Basic Resource Manager}
Basic Manager 面向那些不需要特殊拓扑建模、也不需要复杂上下文切换机制的资源。它实现统一调度器所需的基础能力，例如资源申请、释放、容量可见性与排队协调，是整个系统的最小实现单元。

\subsection{基于 AOE 的 CPU Manager}
CPU 资源管理器主要面向 AI 编程等使用环境执行、命令运行和测试计算的场景。

\parabf{CPU Manager 中的 Breakdown。}
对 CPU 场景而言，Breakdown 的关键是：在动作结束后尽快释放 CPU 资源，同时保留环境状态，以便下一个动作到来时能够快速恢复，而不必重新构建整个执行环境。

\parabf{CPU Manager 中的 Pool。}
在 Pool 方面，CPU Manager 通过感知节点内资源与任务需求，把相同类型的 CPU 资源池化给不同动作共享使用，并减少资源碎片，提升并发执行效率。对可扩展动作而言，它还能提高并行度以缩短执行时间。

\subsection{基于 EOE 的 GPU Manager}
GPU 资源管理器主要面向奖励模型、多模型服务和教师模型等依赖 GPU 的外部服务。

\parabf{GPU Manager 中的 Breakdown。}
在 GPU 场景中，Breakdown 更加复杂，因为动作往往依赖模型恢复、缓存状态与张量并行配置。GPU Manager 的目标是在动作结束后回收 GPU 资源，同时尽量缩短后续恢复成本。

\parabf{GPU Manager 中的 Pool。}
在 Pool 方面，GPU Manager 需要同时考虑模型并行度、显存、GPU 数量以及拓扑结构。它通过把资源池化给多个模型或服务共享，减少传统静态部署下的过量预留，并在高并发时依据弹性调度器决策动态调整资源占用。
